% STATUS: written 2026-11-05 (Vukman 1990 read in full)
% Source read in full: Vukman 1990 (pp. 181--189). The section reports the two
% trace theorems for symmetric bi-derivations and the closing conjecture. Maksa's
% two papers (references 4 and 5 of Vukman 1990) are atlas-pending, as is the
% second of them ("On the trace of symmetric bi-derivations", the only example).
\section{Symmetric biadditive maps and their traces}\label{sec:higher-multi:symmetric-maps}

Section~\ref{sec:higher-multi:biderivations} defined the trace of a symmetric
biderivation as the diagonal map $f(x)=D(x,x)$
\cite[p.~245]{Vukman1989Bid}. This section studies the trace on its own
\cite[p.~181]{Vukman1990Two}. Maksa introduced the symmetric biderivation
and the map that sends $x$ to $D(x,x)$ \cite[references~4
and~5]{Vukman1990Two}. The two results of Vukman force the ring to commute from
the behaviour of the trace alone \cite[p.~181]{Vukman1990Two}.

The first result is the biderivation version of the product theorem of
Posner \cite[p.~181]{Vukman1990Two}.

\begin{theorem}[{\cite[Theorem~1, p.~185]{Vukman1990Two}}]\label{thm:higher:twotrace}
Let $R$ be a prime ring of characteristic different from two and three.
Suppose there exist symmetric biderivations $D_1\colon R\times R\to R$ and
$D_2\colon R\times R\to R$, such that $f_1(x)\,f_2(x)=0$, $x\in R$ holds,
where $f_1$ and $f_2$ are the traces of $D_1$ and $D_2$, respectively. In
this case either $D_1=0$ or $D_2=0$.
\end{theorem}

The second weakens the commuting hypothesis of
Theorem~\ref{thm:higher:bid-commuting} a second time
\cite[p.~187]{Vukman1990Two}. Now even one extra commutator is allowed
before the ring is forced to commute \cite[p.~187]{Vukman1990Two}.

\begin{theorem}[{\cite[Theorem~2, p.~187]{Vukman1990Two}}]\label{thm:higher:doublycentral}
Let $R$ be a noncommutative prime ring of characteristic different from two
and three. Suppose there exists a symmetric biderivation $D\colon R\times
R\to R$, such that the mapping $x\mapsto [[f(x),x],x]$, where $f$ stands
for the trace of $D$, is centralizing on $R$. In this case $D=0$.
\end{theorem}

Vukman records the general statement as a conjecture
\cite[pp.~188--189]{Vukman1990Two}. For some $n\geq 1$ the $n$-fold
commutator of the trace with $x$ lies in the center for all $x$, and then
the symmetric biderivation vanishes \cite[pp.~188--189]{Vukman1990Two}. He
marks the case of arbitrary $n$ as hard \cite[p.~189]{Vukman1990Two}.

The general case is open, and this section states it as a problem. % claim-ok: own signpost introducing the moved problem, which carries its own citation

\begin{problem}[{\cite[pp.~188--189]{Vukman1990Two}}]\label{prob:open:vukman}
Let $R$ be a noncommutative prime ring of characteristic not $2$, let $D$ be a
symmetric biderivation of $R$, and let $f$ be its trace.
Suppose that for some $n\ge1$ the $n$-fold commutator $[f(x),x]_n$ lies in the
center for every $x\in R$.
Does $D$ then vanish?
\end{problem}

Vukman proves the statement for $n\le2$ \cite[Theorem~2, p.~187]{Vukman1990Two}
and states the general case as a conjecture, which he expects to be hard
\cite[pp.~188--189]{Vukman1990Two}. The two sections of the chapter join on the
diagonal \cite[p.~490]{Yazarli2017Trideriv}. A vanishing trace of a symmetric
biderivation forces the ring to commute
\cite[p.~187]{Vukman1990Two}. The trace of a permuting
triderivation of Section~\ref{sec:higher-multi:n-derivations} is the
second.

Abdio\u{g}lu and Lee prove for a noncommutative prime ring $R$ with extended centroid $C$ that a bi-additive map $B$ of $R\times R$ into $Q_{ml}(R)$ with $[B(x,y),[x,y]]=0$ for all $x,y$ must have the form $B(x,y)=\lambda[x,y]+\mu(x,y)$ with $\lambda$ in $C$ and $\mu$ bi-additive into $C$, and apply the result to Jordan $\sigma$-biderivations and to commuting maps on Lie ideals~\cite{WlABasicFunctionalIdentityWithAp2017}.

Ashraf and Jamal show for a prime ring $R$ that the stated commutativity identities involving the trace $\delta(x)=\Delta(x,\dots,x)$ of a permuting $n$-additive map force $R$ to be commutative, and give further commutativity criteria from two nonzero permuting $n$-derivations and their traces~\cite{WlTracesOfPermutingNAdditiveMaps2014}.

For an associative prime ring of characteristic different from two, \c{C}even and \c{C}ilo\u{g}lu show that every symmetric Jordan bi-derivation is a symmetric bi-derivation, and that a nonzero symmetric Jordan left bi-derivation forces a 2- and 3-torsion free prime ring to be commutative~\cite{WlOnSymmetricJordanAndJordanLeft2018}.
